% Part: normal-modal-logic
% Chapter: tableaux
% Section: simple-S5

\documentclass[../../../include/open-logic-section]{subfiles}

\begin{document}

\olfileid[fr]{nml}{tab}{s5}

\olsection{\usetoken{P}{tableau} simplifiés pour \Log{S5}}

\Log{S5} est correct et complet relativement à la classe des modèles
universels, c'est-à-dire des modèles dans lesquels tout monde est
accessible depuis chaque monde. Dans ces modèles, la relation
d'accessibilité n'apporte aucune information supplémentaire :
« il existe un monde~$w$ tel que $\mSat{M}{!A}[w]$ » équivaut à
« il existe un tel $w$ accessible depuis~$u$ ». Pour \Log{S5}, nous
pouvons donc travailler avec des modèles universels décrits par un
ensemble de mondes et une valuation~$V$. Chaque classe d'équivalence
d'un cadre de \Log{S5} forme un cadre universel. Lorsqu'on restreint
le modèle à la classe d'un monde, on restreint aussi la valuation,
et chaque monde de la classe garde exactement les mêmes successeurs.
La vérité des !!{formula}s y est donc préservée, par récurrence sur
leur construction.

On peut ainsi simplifier les règles des !!{tableau}s. Dans le cas général, les préfixes sont des suites
d'entiers strictement positifs qui permettent de repérer les mondes
accessibles depuis d'autres : $\sigma.n$ désigne un monde accessible
depuis~$\sigma$. Dans les modèles universels utilisés pour \Log{S5}, tout monde est accessible depuis chaque
monde ; il est donc inutile de coder cette information dans les
préfixes. Nous pouvons prendre de simples entiers strictement
positifs comme préfixes. Les règles simplifiées figurent dans le
\olref{tab:rules-S5}.

\begin{table}
  \begin{center}
    \def\arraystretch{3}\def\fCenter{}
    \begin{tabular}{|c|c|}
    \hline
    \iftag{prvBox}{
      \AxiomC{\sFmla{\True}{\Box !A}[n]}
      \RightLabel{\TRule{\True}{\Box}}
      \UnaryInfC{\sFmla{\True}{!A}[m]}
      \DisplayProof
      &
      \AxiomC{\sFmla{\False}{\Box !A}[n]}
      \RightLabel{\TRule{\False}{\Box}}
      \UnaryInfC{\sFmla{\False}{!A}[m]}
      \DisplayProof\\
      $m$ est utilisé & $m$ est nouveau
      \\[1ex]
      \hline}{}
    \iftag{prvDiamond}{
      \AxiomC{\sFmla{\True}{\Diamond !A}[n]}
      \RightLabel{\TRule{\True}{\Diamond}}
      \UnaryInfC{\sFmla{\True}{!A}[m]}
      \DisplayProof
      &
      \AxiomC{\sFmla{\False}{\Diamond !A}[n]}
      \RightLabel{\TRule{\False}{\Diamond}}
      \UnaryInfC{\sFmla{\False}{!A}[m]}
      \DisplayProof\\
      $m$ est nouveau & $m$ est utilisé
      \\[1ex]
      \hline}{}
    \end{tabular}
  \end{center}
  \caption{Règles simplifiées pour \Log{S5}.}
  \ollabel{tab:rules-S5}
\end{table}

\begin{ex}
  Voici un tableau simplifié et fermé qui démontre $\Log{S5} \Proves
  \Ax{5}$, c'est-à-dire $\Diamond!A \lif \Box\Diamond!A$.
  \begin{oltableau}
    [\pFmla{\False}{\Diamond\formula{A} \lif \Box\Diamond \formula{A}}{1},
      just = \TAss
      [\pFmla{\True}{\Diamond \formula{A}}{1}, just = {\TRule{\False}{\lif}[1]}
        [\pFmla{\False}{\Box\Diamond \formula{A}}{1},
          just = {\TRule{\False}{\lif}[1]}
          [\pFmla{\False}{\Diamond \formula{A}}{2},
            just = {\TRule{\False}{\Box}[3]}
            [\pFmla{\True}{\formula{A}}{3},
              just = {\TRule{\True}{\Diamond}[2]}
                [\pFmla{\False}{\formula{A}}{3},
                  just = {\TRule{\False}{\Diamond}[4]}, close]
            ]
          ]
        ]
      ]
    ]
  \end{oltableau}
  Les préfixes~$2$ et~$3$ sont introduits pour désigner des mondes
  témoins. La dernière règle s'applique au préfixe~$3$, déjà utilisé,
  et y donne l'opérande faux, qui ferme la branche.
\end{ex}

\end{document}
